\subsection{What Should the Model Be Trained Against?}
\label{section:fc_value_aware}

The models described above either record observed outcomes or fit predictive objectives. The Dyna-Q instance in this chapter stores the most recent outcome for each visited state-action pair. The controller-model architecture of \citet{HaSchmidhuber2018} trains a variational autoencoder and a mixture-density recurrent network for observation reconstruction and next-latent prediction. Recurrent state-space models combine reconstruction, reward-prediction, and Kullback-Leibler consistency terms. The deep models are trained for predictive fit, which need not rank model errors by their effect on value estimates. Related work includes value-aware model learning \citep{FarahmandVAML2017}, value-targeted regression \citep{AyoubVTR2020}, value equivalence \citep{GrimmVALEQ2020,GrimmPVE2021}, MuZero \citep{Schrittwieser2020MuZero}, and calibration \citep{VoelckerCalib2025}.

\subsubsection{Value-aware model learning}
\label{section:fc_value_aware_vaml}

\citet{FarahmandVAML2017} define the objective for a value-function class $\mathcal{V}$, a distribution $\mu$ over state-action pairs, an estimated transition kernel $\widehat P$ selected from a model class $\mathcal{M}$, and the true transition kernel $P^\star$. Their population VAML loss is
\[
\mathcal{L}_{\mathrm{VAML}}(\widehat P)
=
\int_{\mathcal{S}\times\mathcal{A}}
\sup_{V\in\mathcal{V}}
\left[
\int_{\mathcal{S}}
V(s')
\left(
P^\star(\mathrm{d}s' \mid s,a)
-
\widehat P(\mathrm{d}s' \mid s,a)
\right)
\right]^2
\,\mathrm{d}\mu(s,a).
\]
Here, $(\widehat P V)(s, a) = \int V(s') \widehat P(\mathrm{d} s' \mid s, a)$ is the expected next-state value under the learned model and $(P^\star V)(s, a)$ is the same quantity under the true kernel. The supremum is taken separately at each state-action pair before the squared discrepancy is averaged under $\mu$. A model that places mass on features of the next observation irrelevant to every value function in $\mathcal{V}$ is not penalized; a model that distorts the value expectation is. \citet{FarahmandIterVAML2018} replace the worst-case value class with the value iterate produced at each step of approximate value iteration. Their analysis bounds how model error and approximate Bellman-update error propagate to the final policy under stated sampling, model-class, boundedness, and concentrability conditions. For the corresponding Lipschitz value-function class and under their smoothness conditions, \citet{AsadiWasserstein2018} show that the pointwise VAML objective is proportional to the squared 1-Wasserstein distance between $\widehat P$ and $P^\star$.

VAML assigns no loss to transition-model errors that leave the expected next-state value of every function in $\mathcal{V}$ unchanged at the state-action pairs weighted by $\mu$. Predictive losses, including the reconstruction and latent-consistency terms used by recurrent state-space models, may instead allocate capacity to details that do not affect those expectations. Whether this affects policy performance depends on the task and model class. The experiments in \citet{FarahmandVAML2017} and \citet{GrimmVALEQ2020} compare selected misspecified or capacity-constrained models; they do not establish that the performance gap grows generally with observation dimensionality.

\subsubsection{Value equivalence and MuZero}
\label{section:fc_value_aware_ve}

\citet{GrimmVALEQ2020} define value equivalence with respect to a policy set $\Pi$ and a function set $\mathcal{V}$. With rewards held fixed, transition kernels $\widehat P$ and $P^\star$ are value-equivalent when $(T^\pi_{\widehat P} V)(s) = (T^\pi_{P^\star} V)(s)$ for all $s$ and all $(\pi, V) \in \Pi \times \mathcal{V}$. Enlarging either $\Pi$ or $\mathcal{V}$ weakly reduces the class of value-equivalent models. If the representable model class contains the true model, using all policies and functions leaves only that model; otherwise the class may be empty. A restricted model class can therefore suffice for the selected policies and functions without representing the full environment, provided that it contains a model with the required Bellman equivalences. \citet{GrimmPVE2021} define proper value equivalence by requiring each selected policy to have the same value function under the model and the environment. They show that an optimal policy computed in any model with this property for all deterministic policies is also optimal in the environment. Under a smoothness assumption and a simplified squared-loss representation of MuZero, they also show that a constant multiple of the expected model loss under the behavior policy's stationary distribution upper-bounds the squared proper-value-equivalence loss for that policy. \citet{AyoubVTR2020} use value-targeted regression to construct confidence sets for optimistic planning in finite-horizon episodic MDPs with a known transition-model family. They prove sublinear regret; for linear-mixture models, the bound is $\widetilde O(d\sqrt{H^3T})$ and does not depend on the number of states or actions.

MuZero learns latent dynamics for use in Monte Carlo tree search. \citet{Schrittwieser2020MuZero} jointly train a representation function, a latent dynamics function, and a prediction function. At each unrolled step, the network predicts reward, value, and policy. The targets are the observed reward, a final outcome or an $n$-step bootstrapped return, and the policy returned by Monte Carlo tree search. There is no reconstruction loss; the latent dynamics are never asked to decode observations. The latent dynamics have no direct target for state or observation prediction. Their gradients come from the reward, value, and policy losses, with the reward loss omitted in the board-game experiments. In the reported experiments, MuZero matched AlphaZero on Go, chess, and shogi without access to the game rules. On the 57-game Atari benchmark, its aggregate mean and median human-normalized scores exceeded those of the comparison agents trained in the large-data setting. \citet{Antonoglou2022StochMuZero} factor the latent transition into an action-conditioned afterstate and a stochastic transition from that afterstate to the next latent state. Their tree search alternates decision nodes with chance nodes. MuZero and Stochastic MuZero train latent models on targets used by search rather than reconstructing observations. Their latent dynamics can support planning without representing the full observation distribution.\footnote{\citet{VoelckerCalib2025} show that sampled value-aware losses are uncalibrated for stochastic model classes. The losses can favor models with too little variance, and MuZero-style updates need not recover the per-state value function even when the stochastic model is correct. Their Calibrated VAML loss subtracts an estimate of the variance of model-generated value predictions, which restores calibration for model learning in expectation. They recommend updating the value function with a separate model-based or model-free temporal-difference loss.}\footnote{\citet{DontiAmosKolter2017} train probabilistic predictors against the downstream decision cost of a stochastic program rather than forecast fit alone. Their gradients pass through the solution of the stochastic program rather than through a Bellman backup.}

\subsection{TD-MPC2 for Continuous Control}
\label{section:fc_tdmpc2}

TD-MPC2 \citep{Hansen2024TDMPC2} combines latent-space model predictive control with a decoder-free latent model and a learned terminal value function. The model predicts latent transitions, rewards, and values, and the planner evaluates short imagined trajectories with a value bootstrap. The paper evaluates one hyperparameter configuration across $104$ continuous-control tasks without per-task tuning. In its multi-task scaling experiment, the same training hyperparameters are used at each model size.

\subsubsection{Components and Training Objectives}
\label{section:fc_tdmpc2_components}

TD-MPC2 has five learned components. An encoder $h$ maps an observation $s_t$ to a normalized latent $z_t = h(s_t)$. A latent dynamics model $d$ predicts the next latent, $\hat z_{t+1} = d(z_t, a_t)$. A reward model $R$ predicts the immediate reward from the latent and action, $\hat r_t = R(z_t, a_t)$. A terminal action-value function $Q$ supplies a bootstrap beyond the planning horizon, $\hat q_t = Q(z_t, a_t)$. A policy prior $p$ proposes actions for the planner and supplies the action used in the temporal-difference target. The encoder, dynamics, reward, and value components minimize a joint objective with latent consistency, reward prediction, and temporal-difference value prediction terms. The latent consistency term is $\| d(z_t, a_t) - \operatorname{sg}(h(s_{t+1})) \|^2$, with a stop-gradient on the target encoding. Reward and value prediction use cross-entropy in a discretized, log-transformed space. The policy prior is optimized separately with a maximum-entropy objective that balances $Q(z_t,p(z_t))$ against a policy-entropy term. TD-MPC2 has no observation reconstruction loss and does not decode observations.

\subsubsection{MPPI planning with value bootstrap}
\label{section:fc_tdmpc2_mppi}

Action selection at each environment step uses Model Predictive Path Integral control, a sampling-based optimizer over Gaussian action sequences \citep{Hansen2024TDMPC2}. The planner maintains a sequence of Gaussian distributions $(\mathcal{N}(\mu_h, \sigma_h^2))_{h=t}^{t+H-1}$ over the next $H$ actions, draws a population of candidate sequences, rolls each one forward through the learned latent dynamics, scores the imagined returns, and refits $(\mu_h, \sigma_h)$ to a softmax-weighted average of the top-scoring sequences. The scoring rule is the finite-horizon return augmented with a terminal value bootstrap,
\[
  \mu^\star, \sigma^\star \;=\; \arg\max_{(\mu,\sigma)} \;\mathbb{E}_{(a_t, \ldots, a_{t+H}) \sim \mathcal{N}(\mu, \sigma^2)} \Big[\, \gamma^{H} Q(z_{t+H}, a_{t+H}) \;+\; \sum_{h=t}^{H-1} \gamma^{h-t} R(z_h, a_h) \,\Big],
\]
where the trajectory expectation is taken under the learned latent dynamics $d$ rolled from the current encoded state $z_t = h(s_t)$. The terminal $Q$ estimates return beyond the finite planning horizon. TD-MPC2 runs the planner for a fixed number of iterations, executes $a_t \sim \mathcal{N}(\mu_t^\star, (\sigma_t^\star)^2)$, and plans again after the next observation. The previous planning solution, shifted by one step, initializes the next call. A fraction of candidate sequences are drawn from the policy prior $p$, which also supplies the action used in the $Q$ bootstrap.

\subsubsection{Empirical Results and Scaling}
\label{section:fc_tdmpc2_empirics}

The paper evaluates TD-MPC2 on $104$ continuous-control tasks from four benchmarks, DeepMind Control, Meta-World, ManiSkill2, and MyoSuite, using one hyperparameter configuration. The tasks include action spaces with up to $39$ dimensions, sparse rewards, multi-object manipulation, musculoskeletal motor control, and dog and humanoid locomotion. The reported aggregate score for TD-MPC2 exceeds those of SAC, the original TD-MPC, and DreamerV3 in each benchmark. TD-MPC2 uses the same hyperparameters for every task, whereas SAC and TD-MPC use task-specific settings in these comparisons.

The scaling experiment trains five multi-task agents on the same $80$-task dataset drawn from DeepMind Control and Meta-World, with model sizes from $1$ million to $317$ million parameters. The normalized score rises from $16.0$ at $1$M parameters to $70.6$ at $317$M. The five reported points are approximately linear in log parameter count over the two and a half orders of magnitude examined, and the largest model shows no visible saturation. The same training hyperparameters are used at every model size, with no learning-rate, batch-size, or capacity-specific schedule retuning. The authors do not fit a scaling law, so the evidence is limited to the tested model sizes and task suites.
